\clearpage\thispagestyle{empty}
\pdfbookmark[0]{Quellen zu Instrumente (20. November 1820)}{instrumentsnov1820-sources}
\begin{center}\textbf{Quellen zu Instrumente (20. November 1820)}\end{center}
\noindent Diese KI-Übertragung ins heutige Deutsch folgt \textit{Instrumente} in den \href{https://gdz.sub.uni-goettingen.de/id/PPN236007467}{\textit{Werken}, Band~VI} (Göttingen, 1874), S.~435. Der Erstdruck erschien unter \textit{Göttingen} in den \textit{Göttingischen gelehrten Anzeigen}, Nr.~187, 20. November 1820, \href{https://images.p.sub.uni-goettingen.de/iiif/image/gdz:PPN319721507_1820:00001879/full/1600,/0/default.jpg}{S.~1865} und \href{https://images.p.sub.uni-goettingen.de/iiif/image/gdz:PPN319721507_1820:00001880/full/1600,/0/default.jpg}{S.~1866}, im Exemplar der Niedersächsischen Staats- und Universitätsbibliothek Göttingen.
\par\smallskip\noindent Eine \href{https://www.duden.de/rechtschreibung/Tertie}{Tertie} ist eine Sechzigstelsekunde. Die Wörter über die Empfindlichkeit des Pendels bleiben unter den Strichen in der Werkausgabe lesbar und stehen auch im Erstdruck. KI-Übertragung und KI-Satz.
\clearpage
